\clearpage
\fancyhead[L]{\small HARDY’S PENDELUHR.}
\hypertarget{native-458-hardy}{}
\begin{center}{\large HARDY’S PENDELUHR.}\end{center}

\begin{center}
\rule{0.55\linewidth}{0.4pt}\\[2pt]
Göttingische gelehrte Anzeigen. Stück 101. Seite 1001\,..\,1002. 1826. Juni 26.\\[2pt]
\rule{0.55\linewidth}{0.4pt}
\end{center}

Se. Königl. Hoheit der Herzog von \textsc{Sussex} hat abermals der Universität einen Beweis
Seines gnädigen Wohlwollens gegeben, durch ein kostbares Geschenk an die Sternwarte.
Es besteht in einer astronomischen Pendeluhr von \textsc{Hardy}. Schon seit einigen Monaten
dient sie als Hauptuhr bei den Meridianbeobachtungen, und die Erwartung, welche der
Name des Künstlers und die Vollkommenheit der Arbeit erregten, hat sich dadurch
vollkommen bestätigt. Wir fügen hier die Uebersicht ihres bisherigen Ganges bei, dessen
Gleichförmigkeit nichts zu wünschen übrig lässt. Die Compensation wird auf
\SourceLeaf{459}die allereinfachste Art bewirkt, indem die Pendelstange an ihrem untern Ende anstatt der
sonst gewöhnlichen Linse einen Rahmen trägt, in den ein mit Quecksilber beinahe
angefülltes gläsernes cylindrisches Gefäss von $6\frac{1}{2}$ Zoll Höhe und 2 Zoll Weite gestellt
ist. Wird die Compensation zu schwach befunden, so muss etwas Quecksilber
zugeschüttet, im entgegengesetzten Fall etwas herausgenommen werden, und so wird die
Compensation, wenn der Gang der Uhr erst bei den äussersten Temperaturzuständen
beobachtet ist, einer sehr feinen Regulirung fähig. Die folgende Uebersicht gibt zu
erkennen, dass bei zunehmender Wärme der tägliche Gang ganz allmählig um eine halbe
Secunde retradirt,
 und die Compensation also bisher noch etwas zu schwach ist.

\begin{center}
\emph{Stand und Gang des \textup{\textsc{Hardy}}schen Regulators gegen Sternzeit, auf den Mittag reducirt.}
\end{center}

\begin{center}
\begingroup\small\setlength{\tabcolsep}{5pt}
\begin{tabular}{lr|r|r}
 & & Stand & täglicher Gang \\
\hline
1826 & März 9 & $+0^{\mathrm s}83$ & $+0.42$ \\
 & 11 & $1.67$ & $+0.52$ \\
 & 14 & $3.22$ & $+0.45$ \\
 & 19 & $5.46$ & $+0.35$ \\
 & April 20 & $16.55$ & $+0.14$ \\
 & 21 & $16.69$ & $+0.07$ \\
 & 22 & $16.76$ & $+0.21$ \\
 & Mai 2 & $19.09$ & $+0.16$ \\
 & 8 & $20.05$ & $+0.18$ \\
 & 9 & $20.23$ & $+0.25$ \\
 & 11 & $20.74$ & $+0.01$ \\
\end{tabular}
\par\medskip
\begin{tabular}{lr|r|r}
 & & Stand & täglicher Gang \\
\hline
1826 & Mai 14 & $+20^{\mathrm s}78$ & $+0.08$ \\
 & 16 & $20.94$ & $+0.21$ \\
 & 18 & $21.37$ & $+0.15$ \\
 & 20 & $21.68$ & $+0.02$ \\
 & 23 & $21.73$ & $+0.10$ \\
 & 26 & $22.03$ & $-0.06$ \\
 & 30 & $21.78$ & $-0.18$ \\
 & Juni 11 & $19.57$ & $-0.22$ \\
 & 16 & $18.45$ & $-0.18$ \\
 & 17 & $18.27$ &  \\
\end{tabular}
\endgroup
\end{center}

\clearpage
\fancyhead[L]{\small GAUSS AN BODE.}
\hypertarget{native-459-bode}{}
\begin{center}{\large GAUSS AN BODE.}\end{center}

\begin{center}
\rule{0.55\linewidth}{0.4pt}\\[2pt]
Astronomisches Jahrbuch für 1829. Seite 144\,..\,145. Berlin 1826.\\[2pt]
\rule{0.55\linewidth}{0.4pt}
\end{center}

\begin{flushright}
Göttingen, den 10. Juli 1826.
\end{flushright}

\begingroup\rightskip=0pt plus 1em\relax
Hieneben übersende ich Ihnen, mein hochverehrter Freund, meine Meridian-Beobachtungen
der Pallas und Ceres um die Zeit ihrer diesjährigen Oppositionen zu beliebigem Gebrauch
für Ihr Jahrbuch.
\par\endgroup

\begin{center}
\emph{Beobachtungen der Pallas} 1826.
\end{center}

\begin{center}
\begingroup\small\setlength{\tabcolsep}{4pt}
\begin{tabular}{lr|c|c|c}
 & & \shortstack{M. Z. in\\Göttingen} & \shortstack{Gerade\\Aufsteig.} & \shortstack{Abweich.\\nördl.} \\
\hline
Juni & 21 & $12^{\mathrm u}11^{\mathrm m}35^{\mathrm s}4$ & $272^\circ 31' 4''9$ & $23^\circ 52' 19''4$ \\
 & 22 & $12\ \ 6\ \ 48.7$ & $272\ \ 18\ \ 21.3$ & $23\ \ 51\ \ 24.1$ \\
 & 23 & $12\ \ 2\ \ \ 2.9$ & $272\ \ \ 5\ \ 50.8:$ & $23\ \ 50\ \ \ 1.9$ \\
 & 24 & $11\ 57\ 16.2$ & $271\ 53\ \ 7.0$ & $23\ 48\ 32.4$ \\
 & 25 & $11\ 52\ 30.0$ & $271\ 40\ 30.6$ & $23\ 46\ 39.4$ \\
 & 26 & $11\ 47\ 43.6$ & $271\ 27\ 51.1$ & $23\ 44\ 24.3$ \\
 & 27 & $11\ 42\ 58.0$ & $271\ 15\ 23.9:$ & ---\quad---\quad--- \\
 & 28 & $11\ 38\ 12.4$ & $271\ \ 2\ 55.8$ & $23\ 39\ 10.5$ \\
\end{tabular}
\endgroup
\end{center}

Das schwache Licht des Planeten erschwerte die Beobachtungen und verminderte ihre
Genauigkeit. Am 27. wurde ich durch einen nur eine Minute von der Pallas entfernt
stehenden Stern 10ter Grösse
 irre, und darüber wurde die Zenithdistanz nicht auf den bedeutend schwächern Planeten
eingestellt, auch letzterer nur an zwei Fäden unzuverlässig beobachtet.

\Needspace{15\baselineskip}
\SourceLeaf{460}

\begin{center}
\emph{Beobachtungen der Ceres} 1826.
\end{center}

\begin{center}
\begingroup\small\setlength{\tabcolsep}{4pt}
\begin{tabular}{lr|c|c|c}
 & & \shortstack{M. Z. in\\Göttingen} & \shortstack{Gerade\\Aufsteig.} & \shortstack{Abweich.\\südl.} \\
\hline
Juni & 24 & $12^{\mathrm u}30^{\mathrm m}16^{\mathrm s}6$ & $280^\circ 9' 33''0:$ & $27^\circ 40' 54''6$ \\
 & 25 & $12\ 25\ 22.6$ & $279\ 54\ 59.5$ & $27\ 45\ 11.6$ \\
 & 26 & $12\ 20\ 29.0$ & $279\ 40\ 32.7$ & $27\ 49\ 20.4$ \\
 & 27 & $12\ 15\ 35.2$ & $279\ 26\ \ 1.5$ & $27\ 53\ 25.9$ \\
 & 28 & $12\ 10\ 41.1$ & $279\ 11\ 26.4$ & $27\ 57\ 26.8$ \\
 & 30 & $12\ \ 0\ 52.3$ & $278\ 42\ \ 7.6$ & $28\ \ 5\ 20.8$ \\
Juli & \ 2 & $11\ 51\ \ 3.6$ & $278\ 12\ 49.0$ & $28\ 12\ 57.7$ \\
 & \ 3 & $11\ 46\ \ 9.3$ & $277\ 58\ 11.1$ & $28\ 16\ 37.1$ \\
\end{tabular}
\endgroup
\end{center}

Die Ceres hatte reichlich die 8te Grösse, und die Beobachtungen sind sämmtlich gut: blos
die erste Rectascension gründet sich blos auf zwei Fäden, und war dabei die schwächste
Vergrösserung gebraucht.
